% ─────────────────────────────────────────────────────────────────────────────
% WERSJA POLSKA — ROZDZIAŁ 4 — WERSJE ROBOCZE
% Wszystkie wersje mają identyczną strukturę sekcji.
% Różnią się stylem:
%   Wersja A — formalna, akademicka, metodyczna
%   Wersja B — zwięzła i zorientowana na dane, liczby mówią same za siebie
%   Wersja C — narracyjna i analityczna, skupiona na interpretacji wyników
% ─────────────────────────────────────────────────────────────────────────────


% ═════════════════════════════════════════════════════════════════════════════
% WERSJA A — formalna, metodyczna, akademicka
% ═════════════════════════════════════════════════════════════════════════════

\chapter{Badanie eksperymentalne}

Niniejszy rozdział przedstawia konfigurację eksperymentalną, opis środowiska
symulacyjnego, metodologię ewaluacji oraz wyniki uzyskane przez jedenaście
ocenianych algorytmów SLAM.
Eksperymenty przeprowadzono na jednym zbiorze danych w~celu zapewnienia
pełnej porównywalności wyników między algorytmami.

\section{Konfiguracja eksperymentalna}

\subsection{Platforma robotyczna}

Platformą badawczą jest robot mobilny TurtleBot3~Burger --- dwukołowy robot
różnicowy o~niskim punkcie ciężkości, reprezentatywny dla klasy małych robotów
mobilnych przeznaczonych do nawigacji w~środowiskach wewnętrznych.
Robot porusza się z~prędkością liniową do $0{,}22$~m/s.

\subsection{Konfiguracja sensoryczna}

Robot wyposażony jest w~cztery sensory uzupełniające się wzajemnie.
Trójwymiarowy wirujący skaner LiDAR pracuje z~częstotliwością $33{,}6$~Hz
(okres skanu $\approx30$~ms), dostarczając gęstych chmur punktów 3D.
Kamera stereo Leopard~Imaging LI-AR0234CS-NVIDIA~Hawk o~bazie $100$~mm,
rozdzielczości $1920\times1200$ i~polu widzenia $121{,}5°$ rejestruje obrazy
z~częstotliwością $60$~fps; obliczana jest z~niej mapa głębokości o~rozdzielczości
$640\times360$.
Centrala inercyjna IMU, zintegrowana z~modułem kamery, dostarcza pomiarów
z~częstotliwością $\approx62$~Hz.
Odometria kołowa wyznaczana jest z~enkoderów silników.

\subsection{Środowisko programowe}

Wszystkie algorytmy uruchomiono w~środowisku ROS\,2 (dystrybucja Humble)
na tym samym zestawie sprzętowym.
Dane sensoryczne odtwarzano z~nagrania~(\textit{rosbag}) z~prędkością
rzeczywistą, zapewniając identyczne warunki temporalne dla każdego algorytmu.
Preprocessing obejmował rzutowanie chmury punktów 3D na skan 2D (dla algorytmów
GMapping, Cartographer i~SLAM~Toolbox) oraz przypisanie indeksów pierścieni
LiDAR wymaganych przez LIO-SAM.

\section{Środowisko symulacyjne}

\subsection{Opis środowiska}

Dane zarejestrowano w~symulowanym środowisku magazynowym w~NVIDIA~Isaac~Sim.
Środowisko odwzorowuje typowy magazyn z~regałami ustawionymi w~regularnych rzędach,
co stanowi wyzwanie dla algorytmów opartych na zamykaniu pętli ze względu na
repetytywność geometryczną sceny.
Trasa przejazdu obejmuje pełną pętlę przez wszystkie korytarze między regałami.

\subsection{Rejestracja zbioru danych}

Zbiór danych (bag\_6) zarejestrowano z~wszystkimi tematami sensorycznymi:
chmurą punktów, lewym i~prawym obrazem kamery, mapą głębokości, danymi IMU
i~odometrią kołową.
Synchronizacja czasowa zapewniona jest przez znaczniki czasowe ROS\,2.

\subsection{Ekstrakcja danych referencyjnych}

Symulator NVIDIA~Isaac~Sim dostarcza dokładnych pozycji referencyjnych
(\textit{ground truth}) z~rozdzielczością odpowiadającą częstotliwości
odometrii.
Pozycje referencyjne zapisano jako temat ROS\,2 i~użyto bezpośrednio do obliczenia
metryk ATE i~RPE po wyrównaniu Umeyamy.

\section{Metodologia ewaluacji}

\subsection{Dokładność trajektorii}

Trajektorię estymowaną wyrównuje się do referencji metodą Umeyamy,
a~następnie oblicza się ATE~RMSE i~ATE~mean jako miary błędu globalnego
oraz RPE~RMSE i~RPE~mean dla okien lokalnych.
Wszystkie wartości wyrażone są w~metrach.

\subsection{Jakość mapy}

Dla algorytmów generujących gęstą chmurę punktów 3D wyznaczana jest
odległość Chamfera (CD) między mapą estymowaną a~mapą referencyjną
zbudowaną ze skanów referencyjnych.
Uzupełnieniem jest miara F-score dla progów $5$, $10$ i~$20$~cm,
wyrażona jako harmoniczna precyzji i~czułości przy danym progu odległości.

\subsection{Koszt obliczeniowy}

Średnie obciążenie procesora i~szczytowe zużycie pamięci RAM rejestrowano
w~trakcie działania każdego algorytmu na tym samym zestawie sprzętowym.

\section{Wyniki}

% Auto-generated by scripts/collect_metrics.py
% Bag: bag\_6
\begin{table}[H]
  \centering
  \setlength{\tabcolsep}{4pt}
  \caption[Trajectory and map quality metrics N/A bag6]{Trajectory and map quality metrics N/A bag\_6. ATE and RPE in metres. Row shading: \colorbox{tier1}{\phantom{X}}~ATE$<$0.1\,m, \colorbox{tier2}{\phantom{X}}~0.1--2\,m, \colorbox{tier3}{\phantom{X}}~$>$2\,m.}
  \label{tab:metrics_bag6}
  \renewcommand{\arraystretch}{1.2}
  \resizebox{\linewidth}{!}{%
  \begin{tabular}{lSSSSSSSSS}
    \toprule
    \rowcolor{hdr}
    \textbf{Algorithm} & \textbf{ATE} & \textbf{ATE} & \textbf{ATE} & \textbf{RPE} & \textbf{RPE} & \textbf{CD} & \textbf{F@5} & \textbf{F@10} & \textbf{F@20} \\
    \rowcolor{hdr}
    & \textbf{RMSE} & \textbf{mean} & \textbf{max} & \textbf{RMSE} & \textbf{mean} & \textbf{(m)} & \textbf{cm} & \textbf{cm} & \textbf{cm} \\
    \midrule
    \cellcolor{tier1}GMapping & 0.015 & 0.013 & 0.055 & 0.002 & 0.000 & 2.258 & 0.005 & 0.041 & 0.113 \\
    \cellcolor{tier1}DLIO & 0.042 & 0.035 & 0.376 & 0.003 & 0.002 & 1.598 & 0.043 & 0.099 & 0.196 \\
    \cellcolor{tier2}Cartographer 2D & 0.280 & 0.153 & 1.059 & 0.023 & 0.003 & 0.103 & 0.627 & 0.794 & 0.867 \\
    \cellcolor{tier2}KISS-ICP & 0.505 & 0.229 & 3.098 & 0.023 & 0.017 & 0.570 & 0.053 & 0.232 & 0.449 \\
    \cellcolor{tier2}MOLA Odometry & 0.223 & 0.215 & 0.331 & 0.003 & 0.003 & 0.210 & 0.016 & 0.249 & 0.597 \\
    \cellcolor{tier1}RTABMap LiDAR & 0.004 & 0.004 & 0.014 & 0.002 & 0.002 & 0.049 & 0.598 & 0.951 & 0.999 \\
    \cellcolor{tier3}SLAM Toolbox & 4.942 & 4.447 & 7.430 & 0.013 & 0.003 & 2.990 & 0.003 & 0.025 & 0.133 \\
    \cellcolor{tier2}RTABMap RGB-D & 0.102 & 0.052 & 0.690 & 0.007 & 0.000 & {N/A} & {N/A} & {N/A} & {N/A} \\
    \cellcolor{tier3}FAST-LIO2 & 3.974 & 3.485 & 8.558 & 0.028 & 0.007 & 2.756 & 0.003 & 0.046 & 0.133 \\
    \cellcolor{tier3}ORB-SLAM3 Stereo & 2.587 & 2.421 & 4.663 & 0.211 & 0.142 & {N/A} & {N/A} & {N/A} & {N/A} \\
    \cellcolor{tier3}ORB-SLAM3 S+IMU & 163.928 & 145.218 & 344.464 & {N/A} & {N/A} & {N/A} & {N/A} & {N/A} & {N/A} \\
    \cellcolor{tier2}cuVSLAM & 0.712 & 0.656 & 1.345 & 0.007 & 0.006 & 3.293 & 0.003 & 0.024 & 0.061 \\
    \cellcolor{tier1}lidarslam\_ros2 & 0.034 & 0.032 & 0.058 & 0.007 & 0.006 & 0.327 & 0.018 & 0.234 & 0.512 \\
    \cellcolor{tier1}EKF & 0.007 & 0.003 & 0.051 & 0.009 & 0.003 & 0.085 & 0.195 & 0.723 & 0.978 \\
    \bottomrule
  \end{tabular}}%
  \renewcommand{\arraystretch}{1}
\end{table}


\subsection{Dokładność trajektorii}

\begin{figure}[H]
\centering
% Row 1 — sensor fusion + LiDAR odometry
\subcaptionbox{EKF Dead Reckoning\\ ATE = 0.007\,m}
  {\includegraphics[width=0.32\textwidth]{trajectory_ekf}}%
\hfill
\subcaptionbox{KISS-ICP\\ ATE = 0.505\,m}
  {\includegraphics[width=0.32\textwidth]{trajectory_kiss_icp}}%
\hfill
\subcaptionbox{MOLA Odometry\\ ATE = 0.223\,m}
  {\includegraphics[width=0.32\textwidth]{trajectory_mola_odometry}}
\\[0.8em]
% Row 2 — 2D LiDAR SLAM
\subcaptionbox{Cartographer 2D\\ ATE = 0.280\,m}
  {\includegraphics[width=0.32\textwidth]{trajectory_cartographer2d}}%
\hfill
\subcaptionbox{GMapping\\ ATE = 0.015\,m}
  {\includegraphics[width=0.32\textwidth]{trajectory_gmapping}}%
\hfill
\subcaptionbox{SLAM Toolbox\\ ATE = 4.942\,m}
  {\includegraphics[width=0.32\textwidth]{trajectory_slam_toolbox}}
\\[0.8em]
% Row 3 — 3D LiDAR SLAM + LiDAR-IMU
\subcaptionbox{RTABMap LiDAR\\ ATE = 0.004\,m}
  {\includegraphics[width=0.32\textwidth]{trajectory_rtabmap}}%
\hfill
\subcaptionbox{FAST-LIO2\\ ATE = 3.974\,m}
  {\includegraphics[width=0.32\textwidth]{trajectory_fast_lio}}%
\hfill
\subcaptionbox{DLIO\\ ATE = 0.042\,m}
  {\includegraphics[width=0.32\textwidth]{trajectory_dlio}}
\\[0.8em]
% Row 4 — LiDAR-IMU SLAM + visual
\subcaptionbox{lidarslam\_ros2\\ ATE = 0.034\,m}
  {\includegraphics[width=0.32\textwidth]{trajectory_lidarslam}}%
\hfill
\subcaptionbox{cuVSLAM\\ ATE = 0.712\,m}
  {\includegraphics[width=0.32\textwidth]{trajectory_vslam_bag6}}%
\hfill
\subcaptionbox{RTABMap RGB-D\\ ATE = 0.102\,m}
  {\includegraphics[width=0.32\textwidth]{trajectory_rtabmap_rgbd}}
\\[0.8em]
% Row 5 — ORB-SLAM3 (failure cases)
\subcaptionbox{ORB-SLAM3 Stereo\\ ATE = 2.587\,m}
  {\includegraphics[width=0.32\textwidth]{trajectory_orbslam_stereo}}%
\hfill
\subcaptionbox{ORB-SLAM3 Stereo+IMU\\ ATE = 163.928\,m}
  {\includegraphics[width=0.32\textwidth]{trajectory_orbslam_stereo_imu}}%
\hfill
\begin{minipage}[b]{0.32\textwidth}\end{minipage}% empty slot
\caption{Estimated trajectories for all evaluated algorithms (colour)
  against the ground-truth reference (dashed black line).
  All trajectories are Umeyama-aligned to the ground truth.
  Circle: start, square: end.}
\label{fig:trajectories_bag6}
\end{figure}


Najniższy ATE~RMSE uzyskał RTABMap~LiDAR ($0{,}004$~m), co plasuje go znacząco
powyżej wszystkich pozostałych algorytmów.
W~grupie wyników dobrych ($\text{ATE} < 0{,}05$~m) znalazły się ponadto EKF
($0{,}007$~m), GMapping ($0{,}015$~m), lidarslam\_ros2 ($0{,}034$~m)
i~DLIO ($0{,}042$~m).
Algorytmy Cartographer~2D ($0{,}280$~m), MOLA~Odometry ($0{,}223$~m),
KISS-ICP ($0{,}505$~m), RTABMap~RGB-D ($0{,}102$~m) i~cuVSLAM ($0{,}712$~m)
osiągnęły wyniki pośrednie.
Trzy algorytmy wykazały znaczący błąd: SLAM~Toolbox ($4{,}942$~m),
FAST-LIO2 ($3{,}974$~m) i~ORB-SLAM3~Stereo ($2{,}587$~m);
wersja ORB-SLAM3 z~IMU uległa katastrofalnemu awarii (ATE $163{,}9$~m).

\subsection{Jakość mapy}

RTABMap~LiDAR uzyskał najniższą odległość Chamfera ($0{,}049$~m)
i~najwyższe F-score we wszystkich progach (F@20\,cm $= 0{,}999$),
potwierdzając dominację zarówno w~dokładności trajektorii, jak i~jakości mapy.
EKF ($\text{CD} = 0{,}085$~m) i~Cartographer~2D ($0{,}103$~m) uzyskały
podobną jakość geometryczną mapy, choć produkują ją odmiennymi metodami.
Algorytmy 2D~SLAM (GMapping, SLAM~Toolbox) wykazały niskie wartości F-score
mimo akumulacji skanów 3D, co wynika z~płaskiej struktury ich map.
Algorytmy wizyjne i~te, które nie ukończyły trasy, uniemożliwiły obliczenie
metryk 3D.

\subsection{Koszt obliczeniowy}

% TODO: uzupełnić danymi CPU/RAM z eksperymentów

\section{Dyskusja}

\subsection{Metody oparte na LiDAR}

RTABMap~LiDAR wykazał najwyższą dokładność spośród wszystkich ocenianych
algorytmów, co zawdzięcza mechanizmowi zamykania pętli opartemu na deskryptorach
geometrycznych i~globalnej optymalizacji grafu.
GMapping osiągnął zaskakująco niski ATE ($0{,}015$~m), co wskazuje na skuteczne
działanie filtra Rao-Blackwellized w~regularnym środowisku magazynowym.
SLAM~Toolbox wykazał natomiast poważny błąd ($4{,}942$~m) spowodowany fałszywymi
zamknięciami pętli wynikającymi z~repetytywności geometrii regałów.

\subsection{Metody LiDAR-IMU}

FAST-LIO2 nie ukończył skutecznie trasy (ATE $3{,}974$~m), co wskazuje na trudności
z~inicjalizacją lub utratę śledzenia w~warunkach magazynowych.
Wynik lidarslam\_ros2 ($0{,}034$~m) potwierdza skuteczność ciasnego sprzężenia
LiDAR-IMU przy odpowiedniej kalibracji.

\subsection{Metody wizyjne}

cuVSLAM ($0{,}712$~m) i~RTABMap~RGB-D ($0{,}102$~m) wykazały porównywalną
skuteczność w~warunkach dobrego oświetlenia.
ORB-SLAM3 nie ukończył trasy w~żadnej konfiguracji; prawdopodobną przyczyną
jest niedopasowanie modelu optycznego do szerokokątnego obiektywu kamery Hawk
o~polu widzenia $121{,}5°$.

\subsection{Przypadki awarii}

Trzy algorytmy --- SLAM~Toolbox, FAST-LIO2 i~ORB-SLAM3 --- nie ukończyły
skutecznie trajektorii.
Wspólną cechą wszystkich awarii jest wrażliwość na repetytywność środowiska
lub niedopasowanie modelu algorytmu do charakterystyk stosowanych sensorów.
Wyniki te stanowią cenne uzupełnienie porównania, wskazując granice
stosowalności poszczególnych podejść.


% ═════════════════════════════════════════════════════════════════════════════
% WERSJA B — zwięzła, zorientowana na dane
% ═════════════════════════════════════════════════════════════════════════════

\chapter{Badanie eksperymentalne}

Rozdział przedstawia konfigurację sprzętową i~programową, opis zbioru danych,
metodologię pomiarów oraz wyniki jedenastu algorytmów SLAM.

\section{Konfiguracja eksperymentalna}

\subsection{Platforma robotyczna}

Robot TurtleBot3~Burger to dwukołowa platforma różnicowa przeznaczona do
nawigacji w~środowiskach wewnętrznych, o~maksymalnej prędkości $0{,}22$~m/s.

\subsection{Konfiguracja sensoryczna}

Platforma wyposażona jest w:
\begin{itemize}
    \item trójwymiarowy wirujący LiDAR, $33{,}6$~Hz, okres skanu $\approx30$~ms;
    \item kamerę stereo NVIDIA~Hawk: baza $100$~mm, $1920\times1200$~px,
          FOV $121{,}5°$, $60$~fps; głębokość $640\times360$~px;
    \item IMU zintegrowane z~modułem kamery, $\approx62$~Hz;
    \item odometrię kołową z~enkoderów silników.
\end{itemize}

\subsection{Środowisko programowe}

Wszystkie algorytmy uruchomiono w~ROS\,2 Humble.
Dane odtwarzano z~nagrania z~prędkością rzeczywistą.
Dla algorytmów 2D zastosowano rzutowanie chmury 3D na skan planarny;
dla LIO-SAM --- przypisanie indeksów pierścieni LiDAR z~kąta elewacji.

\section{Środowisko symulacyjne}

\subsection{Opis środowiska}

Zbiór danych zarejestrowano w~symulowanym magazynie w~NVIDIA~Isaac~Sim,
z~regularnymi rzędami regałów stanowiącymi wyzwanie dla algorytmów
zamknięcia pętli.

\subsection{Rejestracja zbioru danych}

Nagranie bag\_6 obejmuje chmurę punktów, obrazy stereo, głębokość, IMU
i~odometrię kołową zsynchronizowane znacznikami czasowymi ROS\,2.

\subsection{Ekstrakcja danych referencyjnych}

Pozycje referencyjne (\textit{ground truth}) pobrano bezpośrednio z~symulatora
i~użyto do obliczenia metryk po wyrównaniu metodą Umeyamy.

\section{Metodologia ewaluacji}

\subsection{Dokładność trajektorii}

Metryki ATE~RMSE, ATE~mean, RPE~RMSE i~RPE~mean wyznaczane są po wyrównaniu
Umeyamy; wartości w~metrach.

\subsection{Jakość mapy}

Odległość Chamfera (CD) i~F-score przy progach $5$, $10$, $20$~cm porównują
gęstą chmurę punktów mapy z~referencją geometryczną.
Metryka nie dotyczy algorytmów 2D ani tych, które nie ukończyły trasy.

\subsection{Koszt obliczeniowy}

Rejestrowano średnie obciążenie CPU i~szczytowe zużycie RAM na tym samym
zestawie sprzętowym.

\section{Wyniki}

% Auto-generated by scripts/collect_metrics.py
% Bag: bag\_6
\begin{table}[H]
  \centering
  \setlength{\tabcolsep}{4pt}
  \caption[Trajectory and map quality metrics N/A bag6]{Trajectory and map quality metrics N/A bag\_6. ATE and RPE in metres. Row shading: \colorbox{tier1}{\phantom{X}}~ATE$<$0.1\,m, \colorbox{tier2}{\phantom{X}}~0.1--2\,m, \colorbox{tier3}{\phantom{X}}~$>$2\,m.}
  \label{tab:metrics_bag6}
  \renewcommand{\arraystretch}{1.2}
  \resizebox{\linewidth}{!}{%
  \begin{tabular}{lSSSSSSSSS}
    \toprule
    \rowcolor{hdr}
    \textbf{Algorithm} & \textbf{ATE} & \textbf{ATE} & \textbf{ATE} & \textbf{RPE} & \textbf{RPE} & \textbf{CD} & \textbf{F@5} & \textbf{F@10} & \textbf{F@20} \\
    \rowcolor{hdr}
    & \textbf{RMSE} & \textbf{mean} & \textbf{max} & \textbf{RMSE} & \textbf{mean} & \textbf{(m)} & \textbf{cm} & \textbf{cm} & \textbf{cm} \\
    \midrule
    \cellcolor{tier1}GMapping & 0.015 & 0.013 & 0.055 & 0.002 & 0.000 & 2.258 & 0.005 & 0.041 & 0.113 \\
    \cellcolor{tier1}DLIO & 0.042 & 0.035 & 0.376 & 0.003 & 0.002 & 1.598 & 0.043 & 0.099 & 0.196 \\
    \cellcolor{tier2}Cartographer 2D & 0.280 & 0.153 & 1.059 & 0.023 & 0.003 & 0.103 & 0.627 & 0.794 & 0.867 \\
    \cellcolor{tier2}KISS-ICP & 0.505 & 0.229 & 3.098 & 0.023 & 0.017 & 0.570 & 0.053 & 0.232 & 0.449 \\
    \cellcolor{tier2}MOLA Odometry & 0.223 & 0.215 & 0.331 & 0.003 & 0.003 & 0.210 & 0.016 & 0.249 & 0.597 \\
    \cellcolor{tier1}RTABMap LiDAR & 0.004 & 0.004 & 0.014 & 0.002 & 0.002 & 0.049 & 0.598 & 0.951 & 0.999 \\
    \cellcolor{tier3}SLAM Toolbox & 4.942 & 4.447 & 7.430 & 0.013 & 0.003 & 2.990 & 0.003 & 0.025 & 0.133 \\
    \cellcolor{tier2}RTABMap RGB-D & 0.102 & 0.052 & 0.690 & 0.007 & 0.000 & {N/A} & {N/A} & {N/A} & {N/A} \\
    \cellcolor{tier3}FAST-LIO2 & 3.974 & 3.485 & 8.558 & 0.028 & 0.007 & 2.756 & 0.003 & 0.046 & 0.133 \\
    \cellcolor{tier3}ORB-SLAM3 Stereo & 2.587 & 2.421 & 4.663 & 0.211 & 0.142 & {N/A} & {N/A} & {N/A} & {N/A} \\
    \cellcolor{tier3}ORB-SLAM3 S+IMU & 163.928 & 145.218 & 344.464 & {N/A} & {N/A} & {N/A} & {N/A} & {N/A} & {N/A} \\
    \cellcolor{tier2}cuVSLAM & 0.712 & 0.656 & 1.345 & 0.007 & 0.006 & 3.293 & 0.003 & 0.024 & 0.061 \\
    \cellcolor{tier1}lidarslam\_ros2 & 0.034 & 0.032 & 0.058 & 0.007 & 0.006 & 0.327 & 0.018 & 0.234 & 0.512 \\
    \cellcolor{tier1}EKF & 0.007 & 0.003 & 0.051 & 0.009 & 0.003 & 0.085 & 0.195 & 0.723 & 0.978 \\
    \bottomrule
  \end{tabular}}%
  \renewcommand{\arraystretch}{1}
\end{table}


\subsection{Dokładność trajektorii}

\begin{figure}[H]
\centering
% Row 1 — sensor fusion + LiDAR odometry
\subcaptionbox{EKF Dead Reckoning\\ ATE = 0.007\,m}
  {\includegraphics[width=0.32\textwidth]{trajectory_ekf}}%
\hfill
\subcaptionbox{KISS-ICP\\ ATE = 0.505\,m}
  {\includegraphics[width=0.32\textwidth]{trajectory_kiss_icp}}%
\hfill
\subcaptionbox{MOLA Odometry\\ ATE = 0.223\,m}
  {\includegraphics[width=0.32\textwidth]{trajectory_mola_odometry}}
\\[0.8em]
% Row 2 — 2D LiDAR SLAM
\subcaptionbox{Cartographer 2D\\ ATE = 0.280\,m}
  {\includegraphics[width=0.32\textwidth]{trajectory_cartographer2d}}%
\hfill
\subcaptionbox{GMapping\\ ATE = 0.015\,m}
  {\includegraphics[width=0.32\textwidth]{trajectory_gmapping}}%
\hfill
\subcaptionbox{SLAM Toolbox\\ ATE = 4.942\,m}
  {\includegraphics[width=0.32\textwidth]{trajectory_slam_toolbox}}
\\[0.8em]
% Row 3 — 3D LiDAR SLAM + LiDAR-IMU
\subcaptionbox{RTABMap LiDAR\\ ATE = 0.004\,m}
  {\includegraphics[width=0.32\textwidth]{trajectory_rtabmap}}%
\hfill
\subcaptionbox{FAST-LIO2\\ ATE = 3.974\,m}
  {\includegraphics[width=0.32\textwidth]{trajectory_fast_lio}}%
\hfill
\subcaptionbox{DLIO\\ ATE = 0.042\,m}
  {\includegraphics[width=0.32\textwidth]{trajectory_dlio}}
\\[0.8em]
% Row 4 — LiDAR-IMU SLAM + visual
\subcaptionbox{lidarslam\_ros2\\ ATE = 0.034\,m}
  {\includegraphics[width=0.32\textwidth]{trajectory_lidarslam}}%
\hfill
\subcaptionbox{cuVSLAM\\ ATE = 0.712\,m}
  {\includegraphics[width=0.32\textwidth]{trajectory_vslam_bag6}}%
\hfill
\subcaptionbox{RTABMap RGB-D\\ ATE = 0.102\,m}
  {\includegraphics[width=0.32\textwidth]{trajectory_rtabmap_rgbd}}
\\[0.8em]
% Row 5 — ORB-SLAM3 (failure cases)
\subcaptionbox{ORB-SLAM3 Stereo\\ ATE = 2.587\,m}
  {\includegraphics[width=0.32\textwidth]{trajectory_orbslam_stereo}}%
\hfill
\subcaptionbox{ORB-SLAM3 Stereo+IMU\\ ATE = 163.928\,m}
  {\includegraphics[width=0.32\textwidth]{trajectory_orbslam_stereo_imu}}%
\hfill
\begin{minipage}[b]{0.32\textwidth}\end{minipage}% empty slot
\caption{Estimated trajectories for all evaluated algorithms (colour)
  against the ground-truth reference (dashed black line).
  All trajectories are Umeyama-aligned to the ground truth.
  Circle: start, square: end.}
\label{fig:trajectories_bag6}
\end{figure}


RTABMap~LiDAR osiągnął ATE~RMSE $0{,}004$~m --- najniższy wynik spośród
wszystkich algorytmów.
EKF ($0{,}007$~m), GMapping ($0{,}015$~m), lidarslam\_ros2 ($0{,}034$~m)
i~DLIO ($0{,}042$~m) zamknęły stawkę dokładnych wyników (ATE poniżej $0{,}05$~m).
Pośrednią dokładność uzyskały RTABMap~RGB-D ($0{,}102$~m), MOLA ($0{,}223$~m),
Cartographer ($0{,}280$~m), KISS-ICP ($0{,}505$~m) i~cuVSLAM ($0{,}712$~m).
SLAM~Toolbox ($4{,}942$~m), FAST-LIO2 ($3{,}974$~m) i~ORB-SLAM3~Stereo
($2{,}587$~m) nie ukończyły skutecznie trasy.
ORB-SLAM3 z~IMU wykazał katastrofalny błąd ($163{,}9$~m).

\subsection{Jakość mapy}

Najlepszą jakość geometryczną mapy uzyskał RTABMap~LiDAR
(CD $0{,}049$~m, F@20\,cm $= 0{,}999$).
EKF (CD $0{,}085$~m) i~Cartographer~2D (CD $0{,}103$~m) zajęły kolejne miejsca.
Algorytmy 2D~SLAM, mimo dobrej dokładności trajektorii, wykazały niskie F-score
ze względu na płaską strukturę budowanych map.

\subsection{Koszt obliczeniowy}

% TODO: uzupełnić danymi CPU/RAM z eksperymentów

\section{Dyskusja}

\subsection{Metody oparte na LiDAR}

RTABMap~LiDAR dominuje w~obu kategoriach dokładności.
GMapping sprawdził się w~regularnym środowisku, ale SLAM~Toolbox doznał awarii
na tych samych danych z~powodu fałszywych zamknięć pętli wynikających
z~powtarzalności geometrii regałów.
Oba algorytmy 2D generują mapy płaskie, co odzwierciedlają niskie F-score.
KISS-ICP i~MOLA~Odometry, jako metody bez zamykania pętli, akumulują dryf
zgodnie z~oczekiwaniami.

\subsection{Metody LiDAR-IMU}

lidarslam\_ros2 (ATE $0{,}034$~m) potwierdził skuteczność ciasnego sprzężenia.
FAST-LIO2 nie ukończył trasy (ATE $3{,}974$~m) --- prawdopodobnie z~powodu
trudności inicjalizacyjnych lub utraty śledzenia w~powtarzalnym środowisku.

\subsection{Metody wizyjne}

RTABMap~RGB-D (ATE $0{,}102$~m) wypadł lepiej niż cuVSLAM ($0{,}712$~m).
ORB-SLAM3 zawiódł w~obu konfiguracjach; szerokokątny obiektyw kamery Hawk
($121{,}5°$ FOV) jest niezgodny z~modelem optycznym ORB-SLAM3.

\subsection{Przypadki awarii}

SLAM~Toolbox, FAST-LIO2 i~ORB-SLAM3 nie ukończyły trasy.
Awarie ujawniają dwie klasy problemów: podatność na repetytywność środowiska
(SLAM~Toolbox) oraz niezgodność modelu sensora z~algorytmem (ORB-SLAM3).


% ═════════════════════════════════════════════════════════════════════════════
% WERSJA C — narracyjna, analityczna, skupiona na interpretacji
% ═════════════════════════════════════════════════════════════════════════════

\chapter{Badanie eksperymentalne}

Ocena algorytmów SLAM wymaga jednolitych warunków eksperymentalnych,
które pozwolą przypisać różnice w~wynikach właściwościom algorytmów,
a~nie zmienności danych lub sprzętu.
Niniejszy rozdział opisuje, jak takie warunki zostały zapewnione,
oraz przedstawia i~interpretuje uzyskane wyniki.

\section{Konfiguracja eksperymentalna}

\subsection{Platforma robotyczna}

Wybór platformy TurtleBot3~Burger był podyktowany jej reprezentatywnością
dla klasy małych robotów mobilnych działających w~środowiskach wewnętrznych.
Robot jest niewielki, porusza się z~małą prędkością (do $0{,}22$~m/s)
i~dysponuje ograniczonymi zasobami mechanicznymi, co odpowiada typowym
zastosowaniom przemysłowym.

\subsection{Konfiguracja sensoryczna}

Bogaty zestaw sensorów pozwolił testować algorytmy ze wszystkich
interesujących rodzin bez konieczności zmiany platformy.
Trójwymiarowy wirujący LiDAR pracujący z~częstotliwością $33{,}6$~Hz dostarcza
gęstych chmur punktów niezależnych od oświetlenia, lecz podatnych na zniekształcenia
geometryczne przy ruchu robota.
Kamera stereo NVIDIA~Hawk o~bazie $100$~mm i~szerokim polu widzenia $121{,}5°$
dostarcza obrazów stereo i~mapy głębokości, jednak jej niestandardowy kąt
widzenia okazał się istotnym czynnikiem wpływającym na wyniki algorytmów
wizyjnych.
Centrala inercyjna IMU, zintegrowana z~kamerą i~pracująca z~częstotliwością
$62$~Hz, umożliwia deskewing skanów i~ciasne sprzężenie z~innymi sensorami.
Odometria kołowa stanowi tanie, choć podatne na błędy uzupełnienie zestawu.

\subsection{Środowisko programowe}

Wspólne środowisko ROS\,2~Humble eliminuje zmienność wynikającą z~różnic
w~interfejsach oprogramowania.
Odtwarzanie danych z~nagrania z~prędkością rzeczywistą gwarantuje identyczne
warunki temporalne dla wszystkich algorytmów.
Niezbędny preprocessing --- rzutowanie chmury 3D na skan 2D
i~przypisanie indeksów pierścieni LiDAR --- wykonano jednolicie przed
uruchomieniem algorytmów, które tego wymagają.

\section{Środowisko symulacyjne}

\subsection{Opis środowiska}

Symulowany magazyn w~NVIDIA~Isaac~Sim stanowi reprezentatywne środowisko
dla zastosowań logistycznych.
Regularne rzędy regałów tworzą powtarzalną geometrię, która jest zarówno
typowa dla rzeczywistych zastosowań, jak i~szczególnie trudna dla algorytmów
detekcji i~zamykania pętli --- sceny wyglądają podobnie niezależnie od miejsca
w~korytarzu.
Ta właściwość środowiska okazała się kluczowa dla zrozumienia niektórych
obserwowanych wyników.

\subsection{Rejestracja zbioru danych}

Nagranie bag\_6 obejmuje pełną pętlę przez wszystkie korytarze magazynu
ze wszystkimi tematami sensorycznymi: chmurą punktów LiDAR, obrazami stereo,
mapą głębokości, danymi IMU i~odometrią kołową.

\subsection{Ekstrakcja danych referencyjnych}

Unikalną zaletą środowiska symulacyjnego jest dostępność dokładnych
pozycji referencyjnych.
Dane z~NVIDIA~Isaac~Sim zostały nagrane synchronicznie z~pomiarami sensorycznymi
i~użyte bezpośrednio jako \textit{ground truth} do obliczenia metryk po
wyrównaniu Umeyamy.

\section{Metodologia ewaluacji}

\subsection{Dokładność trajektorii}

ATE mierzy globalną spójność trajektorii --- wykazuje, jak bardzo
estymowana trajektoria odchyla się od referencji po optymalnym wyrównaniu.
RPE uzupełnia obraz o~lokalną jakość odometryczną, mierząc błąd względny
między parami poz odległych o~stały krok.
Połączenie obu metryk pozwala odróżnić algorytmy z~dobrą odometrią lokalną
ale słabym zamykaniem pętli od tych, które osiągają globalną spójność kosztem
wyższego błędu lokalnego.

\subsection{Jakość mapy}

Odległość Chamfera i~miara F-score oceniają, czy geometria trójwymiarowej
mapy wiernie odwzorowuje rzeczywistą strukturę środowiska.
Metryki te mają sens tylko dla algorytmów produkujących gęste chmury
punktów 3D --- algorytmy 2D i~te, które nie ukończyły trasy,
nie mogą być w~ten sposób porównane.

\subsection{Koszt obliczeniowy}

Pomiary CPU i~RAM w~trakcie działania algorytmu na tym samym zestawie
sprzętowym pozwalają ocenić praktyczną stosowalność algorytmu
na platformach o~ograniczonych zasobach.

\section{Wyniki}

% Auto-generated by scripts/collect_metrics.py
% Bag: bag\_6
\begin{table}[H]
  \centering
  \setlength{\tabcolsep}{4pt}
  \caption[Trajectory and map quality metrics N/A bag6]{Trajectory and map quality metrics N/A bag\_6. ATE and RPE in metres. Row shading: \colorbox{tier1}{\phantom{X}}~ATE$<$0.1\,m, \colorbox{tier2}{\phantom{X}}~0.1--2\,m, \colorbox{tier3}{\phantom{X}}~$>$2\,m.}
  \label{tab:metrics_bag6}
  \renewcommand{\arraystretch}{1.2}
  \resizebox{\linewidth}{!}{%
  \begin{tabular}{lSSSSSSSSS}
    \toprule
    \rowcolor{hdr}
    \textbf{Algorithm} & \textbf{ATE} & \textbf{ATE} & \textbf{ATE} & \textbf{RPE} & \textbf{RPE} & \textbf{CD} & \textbf{F@5} & \textbf{F@10} & \textbf{F@20} \\
    \rowcolor{hdr}
    & \textbf{RMSE} & \textbf{mean} & \textbf{max} & \textbf{RMSE} & \textbf{mean} & \textbf{(m)} & \textbf{cm} & \textbf{cm} & \textbf{cm} \\
    \midrule
    \cellcolor{tier1}GMapping & 0.015 & 0.013 & 0.055 & 0.002 & 0.000 & 2.258 & 0.005 & 0.041 & 0.113 \\
    \cellcolor{tier1}DLIO & 0.042 & 0.035 & 0.376 & 0.003 & 0.002 & 1.598 & 0.043 & 0.099 & 0.196 \\
    \cellcolor{tier2}Cartographer 2D & 0.280 & 0.153 & 1.059 & 0.023 & 0.003 & 0.103 & 0.627 & 0.794 & 0.867 \\
    \cellcolor{tier2}KISS-ICP & 0.505 & 0.229 & 3.098 & 0.023 & 0.017 & 0.570 & 0.053 & 0.232 & 0.449 \\
    \cellcolor{tier2}MOLA Odometry & 0.223 & 0.215 & 0.331 & 0.003 & 0.003 & 0.210 & 0.016 & 0.249 & 0.597 \\
    \cellcolor{tier1}RTABMap LiDAR & 0.004 & 0.004 & 0.014 & 0.002 & 0.002 & 0.049 & 0.598 & 0.951 & 0.999 \\
    \cellcolor{tier3}SLAM Toolbox & 4.942 & 4.447 & 7.430 & 0.013 & 0.003 & 2.990 & 0.003 & 0.025 & 0.133 \\
    \cellcolor{tier2}RTABMap RGB-D & 0.102 & 0.052 & 0.690 & 0.007 & 0.000 & {N/A} & {N/A} & {N/A} & {N/A} \\
    \cellcolor{tier3}FAST-LIO2 & 3.974 & 3.485 & 8.558 & 0.028 & 0.007 & 2.756 & 0.003 & 0.046 & 0.133 \\
    \cellcolor{tier3}ORB-SLAM3 Stereo & 2.587 & 2.421 & 4.663 & 0.211 & 0.142 & {N/A} & {N/A} & {N/A} & {N/A} \\
    \cellcolor{tier3}ORB-SLAM3 S+IMU & 163.928 & 145.218 & 344.464 & {N/A} & {N/A} & {N/A} & {N/A} & {N/A} & {N/A} \\
    \cellcolor{tier2}cuVSLAM & 0.712 & 0.656 & 1.345 & 0.007 & 0.006 & 3.293 & 0.003 & 0.024 & 0.061 \\
    \cellcolor{tier1}lidarslam\_ros2 & 0.034 & 0.032 & 0.058 & 0.007 & 0.006 & 0.327 & 0.018 & 0.234 & 0.512 \\
    \cellcolor{tier1}EKF & 0.007 & 0.003 & 0.051 & 0.009 & 0.003 & 0.085 & 0.195 & 0.723 & 0.978 \\
    \bottomrule
  \end{tabular}}%
  \renewcommand{\arraystretch}{1}
\end{table}


\subsection{Dokładność trajektorii}

\begin{figure}[H]
\centering
% Row 1 — sensor fusion + LiDAR odometry
\subcaptionbox{EKF Dead Reckoning\\ ATE = 0.007\,m}
  {\includegraphics[width=0.32\textwidth]{trajectory_ekf}}%
\hfill
\subcaptionbox{KISS-ICP\\ ATE = 0.505\,m}
  {\includegraphics[width=0.32\textwidth]{trajectory_kiss_icp}}%
\hfill
\subcaptionbox{MOLA Odometry\\ ATE = 0.223\,m}
  {\includegraphics[width=0.32\textwidth]{trajectory_mola_odometry}}
\\[0.8em]
% Row 2 — 2D LiDAR SLAM
\subcaptionbox{Cartographer 2D\\ ATE = 0.280\,m}
  {\includegraphics[width=0.32\textwidth]{trajectory_cartographer2d}}%
\hfill
\subcaptionbox{GMapping\\ ATE = 0.015\,m}
  {\includegraphics[width=0.32\textwidth]{trajectory_gmapping}}%
\hfill
\subcaptionbox{SLAM Toolbox\\ ATE = 4.942\,m}
  {\includegraphics[width=0.32\textwidth]{trajectory_slam_toolbox}}
\\[0.8em]
% Row 3 — 3D LiDAR SLAM + LiDAR-IMU
\subcaptionbox{RTABMap LiDAR\\ ATE = 0.004\,m}
  {\includegraphics[width=0.32\textwidth]{trajectory_rtabmap}}%
\hfill
\subcaptionbox{FAST-LIO2\\ ATE = 3.974\,m}
  {\includegraphics[width=0.32\textwidth]{trajectory_fast_lio}}%
\hfill
\subcaptionbox{DLIO\\ ATE = 0.042\,m}
  {\includegraphics[width=0.32\textwidth]{trajectory_dlio}}
\\[0.8em]
% Row 4 — LiDAR-IMU SLAM + visual
\subcaptionbox{lidarslam\_ros2\\ ATE = 0.034\,m}
  {\includegraphics[width=0.32\textwidth]{trajectory_lidarslam}}%
\hfill
\subcaptionbox{cuVSLAM\\ ATE = 0.712\,m}
  {\includegraphics[width=0.32\textwidth]{trajectory_vslam_bag6}}%
\hfill
\subcaptionbox{RTABMap RGB-D\\ ATE = 0.102\,m}
  {\includegraphics[width=0.32\textwidth]{trajectory_rtabmap_rgbd}}
\\[0.8em]
% Row 5 — ORB-SLAM3 (failure cases)
\subcaptionbox{ORB-SLAM3 Stereo\\ ATE = 2.587\,m}
  {\includegraphics[width=0.32\textwidth]{trajectory_orbslam_stereo}}%
\hfill
\subcaptionbox{ORB-SLAM3 Stereo+IMU\\ ATE = 163.928\,m}
  {\includegraphics[width=0.32\textwidth]{trajectory_orbslam_stereo_imu}}%
\hfill
\begin{minipage}[b]{0.32\textwidth}\end{minipage}% empty slot
\caption{Estimated trajectories for all evaluated algorithms (colour)
  against the ground-truth reference (dashed black line).
  All trajectories are Umeyama-aligned to the ground truth.
  Circle: start, square: end.}
\label{fig:trajectories_bag6}
\end{figure}


Wyniki wyraźnie skupiają się w~trzech grupach.
W~pierwszej, z~ATE poniżej $0{,}05$~m, znalazły się: RTABMap~LiDAR ($0{,}004$~m),
EKF ($0{,}007$~m), GMapping ($0{,}015$~m), lidarslam\_ros2 ($0{,}034$~m)
i~DLIO ($0{,}042$~m).
Wynik EKF jest zaskakujący: algorytm bez geometrycznej korekcji błędów
uzyskał drugie miejsce, co wynika z~krótkiego czasu trwania nagrania,
w~którym dryf odometrii nie zdążył narosnąć do dużych wartości.

Środkową grupę (ATE $0{,}1$--$1{,}0$~m) tworzą: RTABMap~RGB-D ($0{,}102$~m),
MOLA~Odometry ($0{,}223$~m), Cartographer~2D ($0{,}280$~m),
KISS-ICP ($0{,}505$~m) i~cuVSLAM ($0{,}712$~m).
Te algorytmy ukończyły trasę z~błędem nadal akceptowalnym w~wielu zastosowaniach,
lecz wyraźnie ustępują czołówce.

Trzecia grupa obejmuje algorytmy, które nie ukończyły trasy skutecznie:
ORB-SLAM3~Stereo ($2{,}587$~m), FAST-LIO2 ($3{,}974$~m),
SLAM~Toolbox ($4{,}942$~m) i~katastrofalnie ORB-SLAM3 z~IMU ($163{,}9$~m).

\subsection{Jakość mapy}

W~klasyfikacji jakości map trójwymiarowych RTABMap~LiDAR jest bezkonkurencyjny:
odległość Chamfera $0{,}049$~m i~F@20\,cm $= 0{,}999$ oznaczają, że praktycznie
każdy punkt mapy referencyjnej ma odpowiednik w~mapie estymowanej w~odległości
poniżej $20$~cm.
EKF ($0{,}085$~m) i~Cartographer~2D ($0{,}103$~m) zajmują kolejne miejsca,
co jest interesujące: Cartographer buduje wewnętrznie mapę 2D, ale akumulowane
skany 3D mają wysoką jakość geometryczną dzięki dobrej estymacji pozycji.

Wyniki GMapping ilustrują rozbieżność między dokładnością trajektorii
a~jakością mapy 3D: pomimo ATE $0{,}015$~m F-score wynosi jedynie $0{,}113$
przy progu $20$~cm.
Wynika to z~płaskiej struktury mapy 2D budowanej przez GMapping ---
skany 3D są akumulowane bez globalnej spójności w~osi pionowej.

\subsection{Koszt obliczeniowy}

% TODO: uzupełnić danymi CPU/RAM z eksperymentów

\section{Dyskusja}

\subsection{Metody oparte na LiDAR}

RTABMap~LiDAR wykazał, że grafowy SLAM 3D z~geometrycznym zamykaniem pętli
jest w~stanie uzyskać wyjątkową dokładność nawet w~powtarzalnym środowisku
magazynowym.
Zestawienie z~wynikiem SLAM~Toolbox ($4{,}942$~m) --- algorytmu z~tej samej
rodziny 2D~SLAM co GMapping, lecz z~agresywniejszym mechanizmem zamykania pętli
--- uwidacznia krytyczne ryzyko fałszywych zamknięć pętli w~środowiskach
o~regularnej geometrii.
GMapping, używający probabilistycznego filtra cząsteczkowego zamiast
deterministycznego dopasowania cech, okazał się bardziej odporny na tę pułapkę.

\subsection{Metody LiDAR-IMU}

Kontrast między lidarslam\_ros2 ($0{,}034$~m) a~FAST-LIO2 ($3{,}974$~m)
wskazuje, że samo ciasne sprzężenie LiDAR-IMU nie gwarantuje sukcesu.
FAST-LIO2 jest algorytmem iteracyjnego wygładzania stanu (\textit{IEKF})
zaprojektowanym przede wszystkim dla dynamicznych środowisk zewnętrznych
i~mógł nie poradzić sobie z~wolnym ruchem robota w~środowisku wewnętrznym.

\subsection{Metody wizyjne}

Wyniki algorytmów wizyjnych w~dużej mierze determinuje sprzęt kamery.
RTABMap~RGB-D, korzystający z~gotowej mapy głębokości i~wbudowanego zamykania
pętli Bag-of-Words, uzyskał akceptowalny wynik ($0{,}102$~m).
ORB-SLAM3 zawiódł w~obu konfiguracjach; szerokie pole widzenia kamery Hawk
($121{,}5°$) jest niezgodne z~modelem optycznym ORB-SLAM3, który zakłada
nieskrzywiony obiektyw o~wąskim kącie, co powoduje znaczące błędy reprojekcji.

\subsection{Przypadki awarii}

Trzy algorytmy nie ukończyły skutecznie trajektorii.
SLAM~Toolbox ilustruje wrażliwość geometrycznego zamykania pętli na
repetytywność środowiska.
FAST-LIO2 wskazuje na trudności z~dopasowaniem parametrów do powolnego ruchu
w~środowisku wewnętrznym.
ORB-SLAM3 ujawnia problem niezgodności modelu optycznego z~niestandardową kamerą.
Wszystkie trzy przypadki stanowią wartościową informację praktyczną: wybór
algorytmu musi uwzględniać nie tylko średnią dokładność, ale i~warunki,
w~których algorytm całkowicie zawodzi.
